\part{Neural many-body Bloch wavefunctions}

\section{Motivation and scope}

The second calculation tests the neural many-body Bloch construction of
Ref.~\cite{luo} directly in the continuum model benchmarked in Part~I. The
purpose is not to fit coefficients in the one-band ED Hilbert space. Instead,
a graph network dresses continuum Bloch orbitals in coordinate space and the
resulting generalized Slater determinant is optimized against the complete
interacting continuum Hamiltonian. This can generate remote-band content even
though the fixed one-body input is the mean-field first band.

This revision follows Appendix~B, Eqs.~(16)--(23), of Ref.~\cite{luo}.
In particular, the orbital transformation is the direct complex output
\begin{equation}
 J_{\bm k i}(X)=\operatorname{MLP}([\bm k,V_i^t])\in\mathbb C,
 \qquad Q_{\bm k i}=J_{\bm k i}\varphi_{\bm k i}.
 \label{eq:scope-direct-j}
\end{equation}
There is no constant-one skip, residual parameterization, or factor $0.2$;
the implemented transformation is the direct $J_{\bm k i}$ above.
The current source additionally implements three audit corrections: full-$M$
uses an unbiased complex random initialization; the node aggregation is the
literal $\sum_jm_{ij}$ in Eq.~(19), including $j=i$; and the random auxiliary
features $h_i^0,h_{ij}^0$ are explicitly retained in the trainable parameter
tree. Version 3 further replaces the v2 single selected-momentum no-$M$
determinant by a trainable $M_S$ matrix that linearly combines every selected-
momentum Slater determinant in the exact total-$\Gamma$ sector. Earlier one-hot-$M$, no-self-message and v2
single-minor checkpoints are excluded from the v3 comparison.

Both v3 no-$M$ runs and all six v3 outer-$C_3$ branches are complete.  The
internal-equivariance construction investigated in v2 is omitted from the
present results because its threefold evaluation was expensive and did not
improve the controlled energy comparison. Quoted training
energies use the final ten persistent-chain measurements; the corresponding
RMS is a stability diagnostic rather than an independent error bar.


\section{Configuration, fixed Bloch basis, and conventions}

A walker is a first-quantized configuration
\begin{equation}
 X=\{(\bm r_i,l_i)\}_{i=1}^{N_e},\qquad
 \bm r_i\in[0,1)^2,\quad l_i^{\rm code}\in\{0,1\}.
\end{equation}
The fractional coordinate is converted to a Cartesian coordinate on the
$3\times3$ torus. The layer input used by the network is
\begin{equation}
 s_i=2l_i^{\rm code}-1\in\{-1,+1\},
\end{equation}
which is the $l_i=\pm1$ convention of Ref.~\cite{luo}. Particle labels are
only row labels of the determinant and graph; no ordered occupation string is
used.

The nine mesh momenta and plane-wave orbitals are
\begin{equation}
 \bm k_{m_1m_2}=\frac{m_1}{3}\bm b_1+
 \frac{m_2}{3}\bm b_2,\qquad m_1,m_2=0,1,2,
\end{equation}
\begin{equation}
 \phi^0_{\bm k}(\bm r,l)=
 \sum_{\bm G}u^l_{1\bm k\bm G}
 e^{i(\bm k+\bm G)\cdot\bm r}.
 \label{eq:strictbare}
\end{equation}
The reciprocal set contains the same 61 $\bm G$ vectors per layer used for
the single-particle bands in Fig.~\ref{fig:bands}. The coefficients
$u^l_{1\bm k\bm G}$ are obtained once by diagonalizing the continuum
mean-field Hamiltonian. In JAX they are constants, and in PyTorch they are
registered buffers. They are absent from the trainable parameter tree. Thus
every run uses band 1 as its mean-field orbital source; neither ED amplitudes
nor trainable plane-wave coefficients are used.

\section{Appendix-B graph network}

\subsection{Initial physical and auxiliary features}

The three reciprocal harmonics used at each vertex are
$\bm g_a\in\{\bm b_1,\bm b_2,\bm b_1+\bm b_2\}$. The immutable node
feature is
\begin{equation}
 v_i^0=\left[
 \sin(\bm g_a\!\cdot\!\bm r_i),
 \cos(\bm g_a\!\cdot\!\bm r_i),s_i
 \right]\in\mathbb R^7 .
 \label{eq:nodefeature}
\end{equation}
For every ordered pair $(i,j)$, the immutable edge feature is
\begin{equation}
 e_{ij}^0=\left[
 \sin\!\left(\bm k\!\cdot\!(\bm r_i-\bm r_j)\right),
 \cos\!\left(\bm k\!\cdot\!(\bm r_i-\bm r_j)\right),s_i
 \right]\in\mathbb R^{19},
 \label{eq:edgefeature}
\end{equation}
where all nine mesh momenta are included. In particular, the edge layer
channel is the source layer $s_i$, exactly as in Eq.~(16) of
Ref.~\cite{luo}; no additional layer channels are appended.

The complete initial features are
\begin{equation}
 V_i^0=(v_i^0,h_i^0)\in\mathbb R^{32},\qquad
 E_{ij}^0=(e_{ij}^0,h_{ij}^0)\in\mathbb R^{32}.
\end{equation}
The code therefore uses 25 auxiliary node components and 13 auxiliary edge
components. Appendix~B explicitly calls $h_i^0$ and $h_{ij}^0$ ``randomly
initialized parameters.'' They are therefore leaves of the trainable JAX
pytree (and \texttt{nn.Parameter} objects in PyTorch), so natural-gradient
updates act on them together with the other network parameters. The present
code uses one random vector shared over particle labels and one shared over
ordered pairs; sharing preserves permutation equivariance. Each component is
initialized from $\mathcal N(0,1/32)$. The paper does not specify this
distribution or sharing rule, so both remain explicit local conventions.
There is no learned input embedding before Eq.~\eqref{eq:message}.

\subsection{Two message-passing iterations}

For iteration $a=1,2$,
\begin{align}
 Q_{ij}^{a}&=W_Q^{a}E_{ij}^{a-1},\\
 K_{ij}^{a}&=W_K^{a}E_{ij}^{a-1},\\
 m_{ij}^{a}
 &=F^{a}\!\left(\sum_n
 Q_{in}^{a}\otimes K_{nj}^{a}\right)
 \odot G^{a}(E_{ij}^{a-1}),                             \label{eq:message}\\
 h_i^{a}
 &=h_i^{a-1}+U^{a}\!\left(
 [\sum_jm_{ij}^{a},V_i^{a-1}]\right),\\
 h_{ij}^{a}
 &=h_{ij}^{a-1}+H^{a}\!\left(
 [m_{ij}^{a},E_{ij}^{a-1}]\right),\\
 V_i^{a}&=(v_i^0,h_i^{a}),\\
 E_{ij}^{a}&=(e_{ij}^0,h_{ij}^{a}).
 \label{eq:graphupdate}
\end{align}
The outer product in Eq.~\eqref{eq:message} is flattened before $F^a$.
$F^a,G^a,H^a,U^a$ are one-layer real affine maps followed by exact GELU.
The two iterations have independent parameters. No division by $N_e$, no
neighbor-count normalization, and no LayerNorm is applied. Self edges remain
available to both the edge update and node aggregation. This is the literal
$\sum_j$ printed in Luo--Fu Eq.~(19), rather than the ``all other particles''
neighborhood convention of Ref.~\cite{pescia}. The literal physical features
$v_i^0,e_{ij}^0$ are reinserted after every iteration.

\section{Backflow, orbital transformation, and determinant}

The complex backflow matrix is implemented as a real, bias-free
$32\to4$ linear map. Its first and last two outputs are the real and
imaginary parts of a two-component displacement,
\begin{equation}
 \delta\bm r_i(X)=W V_i^2,\qquad
 \widetilde{\bm r}_i=\bm r_i+\delta\bm r_i .
 \label{eq:strictbackflow}
\end{equation}
There is no $\tanh$, clipping, or external backflow scale. The fixed
mean-field coefficients are evaluated only at this transformed coordinate,
\begin{equation}
 \varphi_{\bm k i}(X)=
 \sum_{\bm G}u^{l_i}_{1\bm k\bm G}
 \exp\!\left[i(\bm k+\bm G)\cdot
 (\bm r_i+\delta\bm r_i(X))\right].
 \label{eq:strictphi}
\end{equation}
Consequently, backflow does not update $u^l_{1\bm k\bm G}$. The factors
$\exp[i(\bm k+\bm G)\cdot\delta\bm r_i]$ are generated only by changing the
coordinate argument.

The neural orbital transformation uses the Cartesian two-component mesh
momentum and the final node feature,
\begin{equation}
 [V_i^2,k_x,k_y]\in\mathbb R^{34}
 \xrightarrow{\ 34\to32\ }\mathrm{GELU}
 \xrightarrow{\ 32\to2\ }(J^{\rm R}_{\bm k i},J^{\rm I}_{\bm k i}),
\end{equation}
\begin{equation}
 J_{\bm k i}=J^{\rm R}_{\bm k i}+iJ^{\rm I}_{\bm k i},\qquad
 Q_{\bm k i}=J_{\bm k i}\varphi_{\bm k i}.
 \label{eq:strictonce}
\end{equation}
The direct complex output is not converted to $1+\alpha\Delta J$.
Equation~\eqref{eq:strictonce} is implemented once: $J$ is not included
inside $\varphi$ and then multiplied again. A unit test that sets
$J_{\bm k i}=1/2$ verifies
$\det Q=2^{-N_e}\det\varphi$ to numerical precision.

For the Luo--Fu full-$M$ state,
\begin{align}
 \Phi_{\gamma i}(X)
 &=\sum_{\bm k}M_{\gamma\bm k}Q_{\bm k i}(X),\\
 \Psi_M(X)&=\det_{\gamma i}\Phi_{\gamma i}(X),
 \qquad M\in\mathbb C^{N_e\times9}.
 \label{eq:fullM}
\end{align}
All graph, $W$, $J$, and complex $M$ parameters are trainable. Only one final
generalized Slater determinant is used. Because $M$ is trainable and dense,
the Cauchy--Binet expansion of $\det(MQ)$ generally contains several
total-momentum sectors.

Version 3 removes the dense single-particle matrix $M$ while retaining a
trainable many-body mixing within the exact total-momentum sector.  Let
$\mathcal C_\Gamma^{(N_e)}$ be the set of canonical $N_e$-element subsets
$I$ of the nine mesh momenta satisfying
$\sum_{\bm k\in I}\bm k=\Gamma$ modulo a primitive reciprocal vector.  The
v3 no-$M$ state is
\begin{equation}
 \begin{gathered}
 \Psi_{\Gamma}^{\rm v3}(X)=
 \sum_{I\in\mathcal C_\Gamma^{(N_e)}}(M_S)_{1I}
 \det_{\bm k\in I,\,i}Q_{\bm k i}(X),\\
 M_S\in\mathbb C^{N_D\times|\mathcal C_\Gamma^{(N_e)}|} .
 \end{gathered}
 \label{eq:nom}
\end{equation}
Both $N_e=3$ and $N_e=6$ have 12 such subsets on the $3\times3$ momentum
mesh.  The determinant-mixing matrix $M_S$ is initialized with a split-real
LeCun-normal convention and optimized by the same stochastic-reconfiguration
update as the graph, backflow and orbital-transformation parameters.  Here
$N_D=1$, so $M_S$ has shape $1\times12$ and contributes 24 real trainable
parameters. Thus v3
fixes total momentum exactly to $\Gamma$ but does not freeze the many-body
momentum content to one bare Slater determinant. The v2 state, retained only
as a historical baseline, kept the single subset that minimized the sum of
bare band-1 energies: $\{(2,0),(2,1),(2,2)\}/3$ for $N_e=3$ and
$\{(1,0),(1,1),(1,2),(2,0),(2,1),(2,2)\}/3$ for $N_e=6$.

The v3 matrix $M_S$ is not the Luo--Fu dense matrix $M$.  A dense $M$
mixes momenta at the one-particle orbital level and its Cauchy--Binet
expansion generally spans several total-momentum sectors.  In contrast,
Eq.~\eqref{eq:nom} uses $M_S$ only at the many-body determinant level and
contains exactly the 12 minors in the $\Gamma$ sector.  Each
minor still uses the fully configuration-dependent $Q_{\bm k i}(X)$, so all
graph, backflow and orbital transformations are recomputed before the linear
combination is formed.

\section{Complete-wavefunction outer $C_3$ projection}
\label{sec:outerc3}

The outer construction applies no internal Reynolds average. In v3 its raw
state is the periodic, exact-total-$\Gamma$ sector mixture
\begin{equation}
 D_\theta^{\rm v3}(X)=
 \sum_{I\in\mathcal C_\Gamma^{(N_e)}}(M_S)_{1I}
 \det_{\bm k\in I,\,i}Q_{\bm k i}(X),
\end{equation}
with an unrestricted Luo--Fu message-passing backflow and $J$ network. Thus
neither an individual displacement nor an individual generalized orbital is
forced to be a $C_3$ eigenobject, while the complete raw state remains in the
translation sector $\bm K_{\rm tot}=\Gamma$.

The physical continuum action includes the layer-dependent Bloch-gauge sewing.
With the code's active-rotation convention, define
\begin{equation}
 (\mathcal C^aD_\theta^{\rm v3})(X)
 =G_a(X)D_\theta^{\rm v3}(C_3^aX),
\end{equation}
where $G_0=1$ and $G_a$ is the accumulated product of the physical layer-gauge
factor at each preceding point on the configuration orbit.  For
$\omega=\exp(2\pi i/3)$, the complete-state projector is
\begin{equation}
 \boxed{\begin{aligned}
 \Psi_{m,\theta}^{\rm v3}(X)&=P_mD_\theta^{\rm v3}(X)\\
 &=\frac13\sum_{a=0}^{2}\omega^{-ma}G_a(X)\\[-2pt]
 &\quad\times D_\theta^{\rm v3}(C_3^aX),\qquad m=0,1,2 .
 \end{aligned}}
 \label{eq:outerc3}
\end{equation}
The three independently optimized sectors have the same 78,124 real
parameters as v3 no $M$; projection adds no trainable parameter. Every raw
branch is a complete graph/determinant evaluation with shared parameters.

Crucially, all stochastic quantities use Eq.~\eqref{eq:outerc3}, not one raw
branch.  Metropolis ratios are $|\Psi_m(X')/\Psi_m(X)|^2$, every kinetic or
tunnelling local-energy term contains the complete ratio
$\Psi_m(X')/\Psi_m(X)$, and SR uses
$\mathcal O_\mu(X)=\partial_{\theta_\mu}\log\Psi_m(X)$.  Setting
\texttt{outer\_c3\_projector=False} recovers the no-$M$ wavefunction exactly.

Before production, 64 random configurations gave relative projection norms
0.5411, 0.5929 and 0.6410 for $m=0,1,2$, so no sector was annihilated.  The
maximum character errors in
$\Psi_m(C_3X)/\Psi_m(X)=\omega^m$ were
$6.13\times10^{-14}$, $6.39\times10^{-14}$ and $5.02\times10^{-14}$.
All idempotency/orthogonality residuals for $P_mP_n=\delta_{mn}P_m$ were below
$7.55\times10^{-15}$, and disabling the projector changed the no-$M$ output
by exactly zero.  These tests isolate the new outer-projection logic from the


\section{Initialization and trainable parameter count}

Every affine neural-network kernel uses JAX LeCunNormal,
$\operatorname{std}=1/\sqrt{\mathrm{fan\,in}}$, with zero biases.
The auxiliary $h^0$ convention was specified above. For full-$M$, the real and
imaginary components are independently initialized with half-scale
LeCunNormal,
\begin{equation}
 \mathbb E[(\Re M_{\gamma k})^2]
 =\mathbb E[(\Im M_{\gamma k})^2]=\frac{1}{2N_k},
 \qquad \mathbb E|M_{\gamma k}|^2=\frac1{N_k}.
\end{equation}
Thus every momentum minor is accessible at step zero and there is no preferred
$\Gamma$ subset. Ref.~\cite{luo} states global LeCunNormal network
initialization but does not separately specify $M$; the complex extension and
variance split above are therefore the present user-requested convention, not
an attributed statement from the paper.

The mean-field orbital source is always band 1. However, because $W$ and
$J$ themselves are LeCunNormal rather than forced to $W=0,J=1$, the step-zero
many-body state is not claimed to equal the undressed bare Slater determinant
pointwise. This distinction is recorded to make the initialization auditable.

\begin{table}[H]
\centering
\caption{Trainable real parameter counts. The v3 no-$M$ and outer states use
the $1\times12$ complex determinant mixer $M_S$.}
\label{tab:strictparams}
\resizebox{\columnwidth}{!}{%
\begin{tabular}{@{}ccccc@{}}
\toprule
filling & $N_e$ & full $M$ & no $M$+$M_S$, $\Gamma$ & outer $P_m$+$M_S$\\
\midrule
$1/3$ & 3 & 78,154 & 78,124 & 78,124\\
$2/3$ & 6 & 78,208 & 78,124 & 78,124\\
\bottomrule
\end{tabular}
}
\end{table}
Relative to v2 no-$M$, $M_S\in\mathbb C^{1\times12}$ adds 24 real parameters.
The external $P_m$ projector shares the same raw parameters among all three
rotated branches and adds no trainable parameter.

\section{Hamiltonian and local-energy estimator}

The VMC objective uses the complete continuum Hamiltonian at
$\theta=3^\circ$, screening parameter $\epsilon=5$, and a $3\times3$ torus.
It is not the one-band or five-band projected ED Hamiltonian. The local energy
contains the kinetic term, layer-dependent moir\'e potential, spatially
varying interlayer tunnelling, and periodic Coulomb interaction,
\begin{equation}
 E_{\rm loc}(X)=\frac{(H\Psi)(X)}{\Psi(X)}
 =E_{\rm kin}+E_{\rm moir\acute e}+E_{\rm tun}+E_{\rm C}.
\end{equation}
The Coulomb sum uses reciprocal-shell cutoff 8, dielectric $\epsilon=5$, and
removes the $\bm q=0$ neutral-background component. Tunnelling is evaluated
by flipping one layer index and multiplying the exact wavefunction ratio.

Cartesian derivatives of the complete complex state use the three-point
central stencil
\begin{align}
 \partial_a\Psi/\Psi&\simeq
 \frac{\Psi(X+h\hat a)-\Psi(X-h\hat a)}{2h\Psi(X)},\\
 \partial_a^2\Psi/\Psi&\simeq
 \frac{\Psi(X+h\hat a)+\Psi(X-h\hat a)-2\Psi(X)}
 {h^2\Psi(X)} .
\end{align}
Here $h=1.5\times10^{-3}a_M$. Periodic wrapping is applied after every
coordinate displacement. A previous $h,h/2,h/4$ audit is not reused as a
validation of the new checkpoints; it must be repeated after convergence.

\subsection{Raw-coordinate and full-wavefunction Hamiltonian audit}

The production JAX path is \texttt{src/jax\_neural\_bloch.py:184--248}; the
independent PyTorch reference is
\texttt{src/continuum\_vmc.py:96--206}.  Both implementations first form
$\bm r_i$ from the input \emph{raw} \texttt{positions}.  Consequently,
\begin{align}
 V_{\rm M}(X)&=-2V\sum_{i,a}\cos(\bm g_a\cdot\bm r_i+\phi_{l_i}),\\
 V_{\rm C}(X)&=\sum_{i<j}\sum_{\bm q\ne0}v_{\bm q}
 \cos[\bm q\cdot(\bm r_i-\bm r_j)]
\end{align}
never see $\delta\bm r_i$.  For tunnelling, $t(\bm r_i)$ is evaluated at
that same raw coordinate, while its numerator is the complete
$\Psi(X_{l_i\to1-l_i})$ obtained by a new \texttt{logpsi} call.  The kinetic
ratios similarly use new complete calls,
\begin{equation}
 R_{i\mu}^{\pm}=\frac{\Psi(X_{\bm r_i\to\bm r_i\pm h\hat\mu})}
 {\Psi(X)},
\end{equation}
so each displaced call rebuilds the graph and recomputes every particle's
$\delta\bm r_j$, $J_{\bm k j}$, the $M Q$ contraction, and the determinant.
No backflow displacement is cached by the Hamiltonian.

As a final energy-zero and convention check, a fixed-state regression sets
$\delta\bm r=0$, $J=1$, and a fixed $M$ selecting a known two-orbital bare
Bloch Slater determinant.  The four local-energy components are then compared
with an independent finite-basis evaluation.  This test uses 4,096 walkers,
120 burn sweeps and 10 measurements and is never used for an optimization
update.
\begin{table}[H]
\centering
\caption{Fixed-state Hamiltonian regression in meV for the complete
many-body state.  Parentheses are Monte Carlo standard errors.}
\label{tab:fixedhamiltonian}
\begin{tabular}{@{}lrrr@{}}
\toprule
term & finite-basis & Monte Carlo & difference\\
\midrule
kinetic & 39.4207 & 39.3467(3343) & $-0.0740$\\
moir\'e & $-38.5168$ & $-38.7739(5347)$ & $-0.2571$\\
tunnelling & $-29.9435$ & $-29.6118(2984)$ & $+0.3317$\\
Coulomb & $-9.6769$ & $-10.1671(3354)$ & $-0.4902$\\
\midrule
total & $-38.7164$ & $-39.2060(3354)$ & $-0.4896$\\
\bottomrule
\end{tabular}
\end{table}
Every component and the total agree within 1.5 quoted standard errors.  The
regression therefore finds no common Hamiltonian normalization, Coulomb
background, tunnelling, or energy-zero mismatch at its present precision.

A targeted regression test in
\texttt{tests/test\_jax\_neural\_bloch.py} instruments the real
\texttt{\_encoded\_graph} function.  One local-energy call produces four
complete graph evaluations: base $X$ with shape $(B,N_e,2)$, all plus and all
minus configurations with shape $(2N_eB,N_e,2)$ each, and all one-particle
layer flips with shape $(N_eB,N_e,2)$.  It checks every displaced coordinate
and every flipped layer against an independently constructed array.  The test
passes, as does the JAX--PyTorch full-local-energy parity test to
$2\times10^{-7}$ meV absolute tolerance.  No Hamiltonian input-routing error
was found in this audit.

\section{MCMC and natural-gradient optimization}

There are 4,128 simultaneous walkers. Initial fractional coordinates are
uniform on the torus and layer indices are uniform in $\{0,1\}$. Each sweep
proposes, for every walker, one random-particle Gaussian displacement with
standard deviation 0.06 in fractional supercell coordinates, followed by one
random-particle layer flip. Both use the Metropolis ratio
$\min[1,|\Psi(X^\prime)/\Psi(X)|^2]$. The present jobs use 300 burn-in sweeps
and two sweeps per update. Acceptance for coordinate and layer moves is
recorded separately.

For real parameters $\theta_a$,
\begin{equation}
 O_a(X)=\partial_{\theta_a}\log\Psi_\theta(X),
\end{equation}
\begin{align}
 S_{ab}&=\operatorname{Re}\left[
 \langle O_a^*O_b\rangle-
 \langle O_a^*\rangle\langle O_b\rangle\right],\\
 F_a&=2\operatorname{Re}\left\langle
 (E_{\rm loc}-\langle E_{\rm loc}\rangle)O_a^*\right\rangle .
\end{align}
The update is
\begin{equation}
 (S+\lambda I)\Delta\theta=F,\qquad
 \theta\leftarrow\theta-\eta\Delta\theta,
 \label{eq:strictsr}
\end{equation}
with ordinary, unpreconditioned CG. JAX JVP/VJP operations apply the metric
matrix-free. The retained runs check at every Krylov step from iteration 20.
For the v3.2 outer runs at both fillings, checks occur at iterations
20,25,30,$\ldots$, and additionally whenever the recursive residual predicts
convergence or the iteration cap is reached. Acceptance always uses the
explicitly recomputed authoritative residual
\begin{equation}
 r_{\rm rel}^{\rm true}=\frac{\lVert F-(S+\lambda I)\Delta\theta\rVert_2}
 {\lVert F\rVert_2}
\end{equation}
with an additional Fisher-vector product. At a check, the solve stops when
$r_{\rm rel}^{\rm true}<0.02$ or at 60 CG iterations. At the cap, an update
is accepted only if $r_{\rm rel}^{\rm true}<0.05$ and
$F^{\mathsf T}\Delta\theta>0$. Otherwise the parameter update is rejected and
the following step uses the recovery values $\lambda=0.03$ and
$\eta=10^{-3}$. After two consecutive stable recovery solves, the protocol
returns to $\lambda=0.01$ and $\eta=2\times10^{-3}$. Every trace stores the
true residual, $F^{\mathsf T}\Delta\theta$, force and direction norms,
acceptances, and whether the update was applied. This is an adaptive truncated
natural gradient, not a claim that the $10^{-3}$ linear-system tolerance has
been reached.

\begin{table}[H]
\centering
\caption{Optimization and numerical settings. Quantities not specified by
Ref.~\cite{luo} are explicitly marked as local choices.}
\label{tab:strictprotocol}
\resizebox{\columnwidth}{!}{%
\begin{tabular}{@{}lll@{}}
\toprule
quantity & current value & provenance\\
\midrule
width / message iterations & 32 / 2, unshared & Ref.~\cite{luo}, Appendix B\\
orbital MLP hidden width & 32, GELU & Ref.~\cite{luo}, Appendix B\\
final determinants & 1 & local controlled choice\\
optimizer / implementation & natural gradient / JAX 0.11.1 & Ref.~\cite{luo}\\
learning rate & $2\times10^{-3}$ & Ref.~\cite{luo}\\
walkers per update & 4,128 & Ref.~\cite{luo}\\
affine-kernel initialization & LeCunNormal & Ref.~\cite{luo}\\
$h^0$ initialization & random, trainable & random: Ref.~\cite{luo}; law/sharing: local\\
full-$M$ initialization & complex LeCunNormal & user choice; paper unspecified\\
node aggregation & $\sum_jm_{ij}$, including self & Ref.~\cite{luo}, Eq.~(19)\\
updates in present experiment & 120 & user-specified finite run\\
paper long-run statement & 1,000 for clusters $>3\times3$ & Ref.~\cite{luo}\\
burn / sweeps per update & 300 / 2 & local choice\\
proposal scale & 0.06 fractional & local choice\\
SR damping & $10^{-2}$ & local choice\\
CG range / true-residual rules & 20--60 / 0.02 stop, 0.05 accept & user-specified\\
wavefunction / SR batch & 258 / 258 & local memory choice\\
local-energy batch & 32 & local memory choice\\
finite-difference step & $1.5\times10^{-3}a_M$ & local choice\\
random seed & run-specific; 3184 for both v3.2 outer fillings & local choice\\
precision & JAX 64 bit & local numerical choice\\
\bottomrule
\end{tabular}%
}
\end{table}

All 4,128 configurations are used in the energy, force, and metric estimate at each
update, with sector-specific importance weights during the v3.2 shared-pool phase.
The printed within-step SEM of private chains treats walker values as an instantaneous
ensemble diagnostic; it is not an independent uncertainty. Checkpoints
at initialization, every 10 updates, the latest completed update, and update
120 store both flattened parameters, optimizer recovery state, sampler RNG,
and walker coordinates/layers. Interrupted jobs resume this complete state;
no validation configuration ever enters an update.

\section{Implementation checks completed before production}

Previously completed architecture, parity, observable, and local-energy checks verify:
\begin{enumerate}
 \item particle exchange changes the determinant sign and no other factor;
 \item $J_{\bm k i}$ is applied exactly once;
 \item fixed Bloch coefficients are absent from trainable parameters;
 \item shifted-coordinate Bloch evaluation agrees with an explicit
 plane-wave calculation;
 \item PyTorch and JAX $\log\Psi$ agree at identical parameters;
 \item complete JAX and PyTorch local energies agree;
 \item no-$M$ is invariant under both primitive translations and contains no
 trainable $M$ array;
 \item the complete-state $C_3$ projector satisfies character, idempotency,
 orthogonality, nonzero-sector-norm and projector-disabled no-$M$ regression
 tests before its three-sector production run;
 \item matrix-free force and Fisher-vector products agree with an explicit
 score-Jacobian calculation;
 \item checkpoint export between JAX and PyTorch is lossless;
 \item full-$M$ is reproducibly random, complex, and not one-hot;
 \item $h_i^0,h_{ij}^0$ are nonzero trainable parameters; and
 \item a one-particle graph receives its nonzero $j=i$ message in Eq.~(19);
 \item the 1-RDM uses full replaced-configuration ratios and a fixed bare
 projector, and its bare-Slater full-basis trace is $N_e$; and
 \item kinetic displacements and tunnelling flips rebuild the complete JAX
 graph, while diagonal Hamiltonian factors use raw coordinates.
\end{enumerate}
A one-update end-to-end smoke test also passed MCMC, local-energy evaluation,
truncated SR, residual recording, parameter update, and checkpoint writing.
These checks validate implementation consistency, not physical convergence.
For v3.2, nine focused numerical tests also passed, including weighted complex
SR/Jacobian parity, the uniform-weight limit, amplitude-invariant mixture
normalization, disconnected-overlap rejection, and sparse true-residual checks.
A shared-pool--fork--private-chain smoke test verified identical initial
parameters, separate accepted updates, identical fork pools and distinct branch RNGs.

\section{Controlled v3.2 sector-mixer experiments}

The historical full-$M$ jobs provide the unrestricted baseline.  The two v3
no-$M$ jobs train $M_S$ for 120 updates with 4,128 persistent walkers.
Version 3.2 uses the same shared-pool protocol for the outer runs at both
$\nu=1/3$ and $2/3$: all three parameter arrays
start identically from seed 3184 but update independently from step one, targeting
\begin{equation}
\mathcal L=\frac13\sum_{m=0}^2
\frac{\langle\Psi_m(\theta_m)|H|\Psi_m(\theta_m)\rangle}
{\langle\Psi_m(\theta_m)|\Psi_m(\theta_m)\rangle}.
\end{equation}
During the first 20 updates, each sector contributes 1,376 persistent walkers
to a common 4,128-configuration pool. Every sector is evaluated on the full
pool. MBAR bridge equations estimate relative normalization constants and
sector importance weights~\cite{mbar}; energy, force and metric centering all
use these weights. This samples a probability mixture, not the coherent
superposition density. The equal factor $1/3$ is canceled between each sector's
force and metric, retaining the original per-sector damping convention.
The weight ESS at step 20 is 2,610/2,828/2,864 at $\nu=1/3$ and
1,875/2,518/2,099 at $\nu=2/3$; it measures weight concentration,
not time-series decorrelation. Finite-sample normalization and MCMC mixing
remain limitations. The exact last update-20 parameter state is retained,
without seed screening or best-checkpoint selection. Each branch receives a
copy of all 4,128 configurations and its own optimizer state, thermalizes for
300 proposal sweeps with a distinct RNG, and performs 100 private updates.
The final-ten training mean determines the displayed sector ordering.

\begin{table}[H]
\centering
\caption{Authoritative local calculations used in v3.2. Both fillings' outer
runs use independent
parameters and shared mixture sampling for updates 1--20.}
\label{tab:adaptive-runs}
\resizebox{\columnwidth}{!}{%
\begin{tabular}{@{}llll@{}}
\toprule
$\nu$ & ansatz & $\bm K_{\rm tot}$ & result-directory suffix\\
\midrule
$1/3$ & full $M$ & mixed & \texttt{gpu\_vmc\_adaptive\_cg120\_seed260904/nu1of3\_full\_m}\\
$1/3$ & no $M$+$M_S$ & $\Gamma$ & \texttt{gpu\_vmc\_adaptive\_cg120\_seed260904/nu1of3\_no\_m\_ms}\\
$1/3$ & outer $P_0$+$M_S$ & $\Gamma$, $m=0$ & \texttt{outer\_c3\_shared\_pool\_nu1of3\_same\_seed3184/p0}\\
$1/3$ & outer $P_1$+$M_S$ & $\Gamma$, $m=1$ & \texttt{outer\_c3\_shared\_pool\_nu1of3\_same\_seed3184/p1}\\
$1/3$ & outer $P_2$+$M_S$ & $\Gamma$, $m=2$ & \texttt{outer\_c3\_shared\_pool\_nu1of3\_same\_seed3184/p2}\\
$2/3$ & outer $P_0$+$M_S$ & $\Gamma$, $m=0$ & \texttt{outer\_c3\_shared\_pool\_same\_seed3184/p0}\\
$2/3$ & outer $P_1$+$M_S$ & $\Gamma$, $m=1$ & \texttt{outer\_c3\_shared\_pool\_same\_seed3184/p1}\\
$2/3$ & outer $P_2$+$M_S$ & $\Gamma$, $m=2$ & \texttt{outer\_c3\_shared\_pool\_same\_seed3184/p2}\\
$2/3$ & full $M$ & mixed & \texttt{gpu\_vmc\_adaptive\_cg120\_seed260904/nu2of3\_full\_m}\\
$2/3$ & no $M$+$M_S$ & $\Gamma$ & \texttt{gpu\_vmc\_adaptive\_cg120\_seed260904/nu2of3\_no\_m\_ms}\\
\bottomrule
\end{tabular}%
}
\end{table}

\subsection{Training-chain diagnostics}

Table~\ref{tab:adaptive-training} summarizes the final ten attempted updates
for each frozen trace.
The mean and RMS describe the persistent training chain and are not independent
energy estimates.  A rejected update contributes its sampled local energy to
the diagnostic trace but leaves the parameters unchanged. All three v3
$\nu=1/3$ branches accepted every update; at $\nu=2/3$, $P_0$ and $P_2$
accepted every update and $P_1$ rejected three.

\begin{table}[H]
\centering
\caption{Final-ten training diagnostics.  Energies are meV per particle;
RMS is training-chain scatter.}
\label{tab:adaptive-training}
\resizebox{\columnwidth}{!}{%
\begin{tabular}{@{}llrrrrrr@{}}
\toprule
$\nu$ & ansatz & $\langle E/N_e\rangle$ & min. $E/N_e$ & RMS &
$\langle r_{\rm rel}\rangle$ & $\langle N_{\rm CG}\rangle$ & reject$_{120}$\\
\midrule
$1/3$ & full $M$ & $-36.2595$ & $-36.3078$ & 0.0286 & 0.0187 & 53.2 & 0\\
$1/3$ & no $M$+$M_S$, $\Gamma$ & $-37.3037$ & $-37.3600$ & 0.0421 & 0.0178 & 34.3 & 0\\
$1/3$ & outer $P_0$+$M_S$ & $-37.3436$ & $-37.3851$ & 0.0343 & 0.0163 & 24.7 & 0\\
$1/3$ & outer $P_1$+$M_S$ & $-32.8446$ & $-32.9533$ & 0.0671 & 0.0173 & 51.3 & 0\\
$1/3$ & outer $P_2$+$M_S$ & $-36.1077$ & $-36.1879$ & 0.0498 & 0.0170 & 27.3 & 0\\
$2/3$ & outer $P_0$+$M_S$ & $-52.5922$ & $-52.7416$ & 0.0867 & 0.0170 & 40.4 & 0\\
$2/3$ & outer $P_1$+$M_S$ & $-51.5740$ & $-51.6848$ & 0.0511 & 0.0168 & 52.8 & 3\\
$2/3$ & outer $P_2$+$M_S$ & $-52.8838$ & $-52.9749$ & 0.0530 & 0.0172 & 38.9 & 0\\
$2/3$ & full $M$ & $-52.1216$ & $-52.1622$ & 0.0376 & 0.0402 & 51.7 & 18\\
$2/3$ & no $M$+$M_S$, $\Gamma$ & $-52.7666$ & $-52.8568$ & 0.0569 & 0.0354 & 54.4 & 16\\
\bottomrule
\end{tabular}%
}
\end{table}
The v3 no-$M$ final-ten position/layer acceptances are 0.672/0.624 and
0.649/0.623 at $\nu=1/3$ and $2/3$.  At $\nu=1/3$, outer $P_0,P_1,P_2$
give 0.669/0.639, 0.699/0.592 and 0.679/0.535, with mean local-energy
variances 55.0, 124.4 and $75.9~\mathrm{meV}^2$.  At $\nu=2/3$, the
$P_0,P_1,P_2$ tails give 0.648/0.592, 0.656/0.613 and 0.650/0.614, with
variances 652.8, 735.1 and $673.4~\mathrm{meV}^2$. Thus a smooth energy trace does not establish a
zero-variance eigenstate.

\subsection{Training-tail comparison}

Following the optimization convention of Ref.~\cite{luo}, the public energy
comparison uses the persistent training ensemble.  For each
trace we average the final ten measurements and quote their RMS fluctuation;
the RMS is a correlated training-chain stability diagnostic, not an
independent standard error.  Among the three outer sectors, $P_0$ is lowest at
$\nu=1/3$, with $\langle E/N_e\rangle_{\rm tail}=-37.3436$ meV and
$\Delta E_{\rm 5b}=+0.0497$ meV. Its $P_1$ and $P_2$ final-ten means remain
0.0586 and 0.0568 meV below the preceding ten-step means.
At $\nu=2/3$, the v3.2 shared-pool run
gives $P_2$ as the lowest sector, with
$\langle E/N_e\rangle_{\rm tail}=-52.8838$ meV and
$\Delta E_{\rm 5b}=-0.1583$ meV. These are correlated optimization-chain
statistics, not independent standard errors. The $P_1$ final-ten mean is still
0.1433 meV below the preceding ten-step mean, so completion of 120 updates
does not establish convergence or seed robustness.

Figure~\ref{fig:sector-summary} retains all three $m$ sectors outside Fig.~\ref{fig:nbresults}.
Its rows correspond to $\nu=1/3$ and $2/3$; columns show optimization traces,
final-ten sector offsets, and separate four-method timing comparisons.
Rejection statistics are retained in Table~\ref{tab:adaptive-training}, not a plot panel.
At $\nu=1/3$, all displayed v3 runs have zero rejected updates. A rejected update leaves the parameters unchanged but
its sampled local energy remains part of the diagnostic trace.
% Preserve the established Fig. 6--7 numbering.
\setcounter{figure}{5}
\begin{figure*}[t]
\centering
\includegraphics[width=\textwidth]{fig6_neural_bloch_results.pdf}
\caption{\textbf{Controlled neural-Bloch optimization, multiband energy,
band mixing and CG stability.} Top row: $\nu=1/3$; bottom row: $\nu=2/3$.
Blue and red denote full $M$ and v3 no $M$+$M_S$; orange denotes the lowest
final-ten v3 outer sector, $P_0$ at $\nu=1/3$ and $P_2$ at $\nu=2/3$. (a,e)
Persistent-walker training energies. (b,f) One- through five-band ED energies
and final-ten training means; shaded widths are the corresponding RMS, not
independent error bars. (c,g) Bare-band populations from the
complete-wavefunction-ratio 1-RDM, including the lowest-training outer sector
and five-band ED in gray. (d,h)
Explicit true CG residuals; horizontal lines mark the 0.02 early-stop and
0.05 cap-acceptance thresholds. The gray line marks update 20. At both
fillings, all 120 updates already optimize a projected state: updates 1--20
use shared-mixture importance weighting and updates 21--120 use its private
4,128-walker chain, retaining the sector's parameters and optimizer state.
Rejected-update counts are
summarized separately in the text.}
\label{fig:nbresults}
\end{figure*}

\begin{figure*}[t]
\centering
\includegraphics[width=\textwidth]{fig6_training_summary.pdf}
\caption{\textbf{Outer-$C_3$ optimization, sector splitting and efficiency.}
(a,d) All three sector-energy traces at $\nu=1/3$ and $2/3$, respectively,
with the complete 120-update history and insets resolving updates 21--120.
The dotted line marks the end of shared mixture sampling at update 20;
projected parameters update independently throughout both fillings' runs.
(b,e) Final-ten energy offsets of the externally projected $P_m$ states,
measured from the lowest sector at each filling. Error bars are RMS
fluctuations of the correlated training tail, not independent standard errors.
At $\nu=1/3$, $m=0$ is selected, while $m=1,2$ lie 4.499 and 1.236 meV per
particle higher. At $\nu=2/3$, $m=2$ is selected, with $m=0,1$ higher by
0.292 and 1.310 meV per particle. Incomplete convergence, particularly in
$P_1$, prevents interpreting these splittings as a converged FCI manifold.
(c,f) Separate nine-cell timing comparisons at $\nu=1/3$ and $2/3$ on a
common logarithmic seconds axis. Five-band ED in $\Gamma$ takes 0.349 s
(median of five fresh processes) and 11.0 min, respectively, for four eigenpairs.
Full $M$ and no $M$+$M_S$ show their individual 120-update job times,
excluding validation, export and observables. For the
hatched outer bars, 7.5 and 26.1 min/sector at $\nu=1/3$ and $2/3$ are
the respective total wall times divided by three candidate states.
This amortized measure is not a standalone
state's runtime. Different scopes and total resource costs are specified in
Table~\ref{tab:ed-qns-efficiency}; no Berry or 27-cell spectrum time enters this comparison.}
\label{fig:sector-summary}
\end{figure*}

Figure~\ref{fig:sector-summary}(c,f) compares the nine-cell QNS runs with
five-band ED restricted to total momentum $\Gamma$, not the complete
Fig.~2--3 production workflow. The measured $\nu=2/3$ ED job took
660.453 s (11.0 min), with a 905,100-dimensional sector, four eigenpairs,
12 configured OpenMP threads and no GPU. It loads a cached five-band
operator and starts from the embedded three-band ground state.
The separately recorded operator setup took 17.129 s and is excluded,
as are the preceding lower-band solves. The 659.286 s internal timer
includes cache loading, the eigensolve, residual and overlap checks;
the 660.453 s process timer additionally includes startup and output.
Neither is a pure eigensolver-kernel timer. CPU process time was not measured.
At $\nu=1/3$, a dedicated five-band $\Gamma$-only timing rerun gives a
median process wall time of 0.349 s across five fresh processes
(range 0.347--0.389 s); median CPU process time is 1.585 s.
The sector dimension is 1,590. It uses the cached interaction, embedded
three-band initial vector and the same four-eigenpair settings as $\nu=2/3$
($\mathrm{tol}=10^{-7}$, $\mathrm{ncv}=64$, maximum 400 iterations).
Median internal and eigensolver-only times are 0.153 s and 0.104 s;
the process timer additionally includes imports, reference checks and output.
All four eigenvalues agree with the saved production results within
$10^{-6}$ meV; existing states are not overwritten. Filesystem caches were
not flushed. This supplementary median is not a contemporaneous repeated
benchmark of the historical $\nu=2/3$ job.
Berry curvature, other momentum sectors and the 27-cell spectrum are excluded.
The full-$M$ and no-$M$+$M_S$ times are reported individually, not as a
two-filling aggregate: 480/327 s at $\nu=1/3$ and 1,467/1,439 s at $\nu=2/3$.
They include startup, initial thermalization and 120 training updates;
optional validation, checkpoint export and observables are excluded.
Full $M$ mixes total momenta and is an unrestricted baseline, whereas no
$M$+$M_S$ and Outer-$C_3$ stay in total $\Gamma$.
The ED calculation returns four eigenpairs, whereas QNS supplies three
finite-training candidate sector states in the full continuum ansatz; the
different accuracy and basis truncation preclude a matched-accuracy speedup claim.
Wall time measures elapsed production time, while accumulated device time
sums simultaneous workers and therefore records resource consumption.  The
ED rows used no GPU; the QNS rows used one A30 and one 16-vCPU affinity per
concurrent job. Software installation, smoke tests, figure rendering and
\TeX{} compilation are excluded. The $\nu=1/3$ outer run took 22.4 minutes,
with 791.39 s user plus 146.84 s system CPU process time (0.261 CPU h).
Its private branches took 6.6, 11.2 and 6.7 minutes for $P_0,P_1,P_2$.
The $\nu=2/3$ outer run took 78.2 minutes including
initialization, thermalization and the shared/private training stages;
its CPU process time was 1,329.31 s user plus 330.77 s system (0.461 CPU h).
Private 100-update branches took 25.6, 36.2 and 29.1 minutes for $P_0,P_1,P_2$.
GPU-hours count lane/worker allocation, not measured kernel execution.
Since the workflow produces three symmetry-sector candidate states, its
amortized cost at $\nu=1/3$ is $22.4/3\simeq7.5$ wall-minutes/sector and
$0.559/3\simeq0.186$ allocated GPU-hours/sector; at $\nu=2/3$ it is
$78.2/3\simeq26.1$ wall-minutes/sector and
$2.059/3=0.686$ allocated GPU-hours/sector. Shared-stage costs are included
before division. The former measures workflow throughput under this schedule,
not the latency of an isolated eigenstate calculation; actual wall time remains
22.4 and 78.2 minutes, respectively. The candidates have not been certified as converged eigenstates.
Separately, on a fixed P2 checkpoint, the sparse-residual SR microbenchmark
reduced the warm solve from 23.94 to 17.58 s (26.6\%) at unchanged acceptance
tolerances; this is not a full-training speedup measurement.

\begin{table}[H]
\centering
\caption{Nine-cell ED/QNS timing. ED is restricted to five bands and $\Gamma$.
GPU-hours are allocation sums; CPU seconds are measured process consumption.
The $\nu=1/3$ ED entry is a five-process median; $\nu=2/3$ is the historical job.
Each amortized outer row divides the complete three-sector training cost by
three; other workflow rows retain their original scopes.}
\label{tab:ed-qns-efficiency}
\resizebox{\columnwidth}{!}{%
\begin{tabular}{@{}llllr@{}}
\toprule
workflow & scope & production resources & wall time & device time\\
\midrule
ED, 5 bands, $\Gamma$ & $\nu=1/3$, four eigenpairs & 12 OpenMP threads & 0.349 s & 1.585 CPU s\\
ED, 5 bands, $\Gamma$ & $\nu=2/3$, four eigenpairs & 12 OpenMP threads & 11.0 min & CPU time not recorded\\
QNS full $M$ training & $\nu=1/3$, one state & A30 & 480 s & 0.133 GPU h\\
QNS no $M$+$M_S$ training & $\nu=1/3$, one state & A30 & 327 s & 0.091 GPU h\\
QNS full $M$ training & $\nu=2/3$, one state & A30 & 1,467 s & 0.408 GPU h\\
QNS no $M$+$M_S$ training & $\nu=2/3$, one state & A30 & 1,439 s & 0.400 GPU h\\
QNS Outer-$C_3$ training & $\nu=1/3$, three sectors & $2\times$A30 & 22.4 min & 0.559 GPU h\\
Outer amortized cost & $\nu=1/3$, per candidate & same schedule & 7.5 min/sector & 0.186 GPU h/sector\\
QNS Outer-$C_3$ training & $\nu=2/3$, three sectors & $2\times$A30 & 78.2 min & 2.059 GPU h\\
Outer amortized cost & $\nu=2/3$, per candidate & same schedule & 26.1 min/sector & 0.686 GPU h/sector\\
Outer observables/export & $\nu=1/3$, three sectors & $2\times$A30 & 21.7 s & 0.009 GPU h\\
Outer observables/export & $\nu=2/3$, three sectors & $2\times$A30 & 25.6 s & 0.011 GPU h\\
\bottomrule
\end{tabular}%
}
\end{table}

\section{Post-training one- and two-body observables}

Figure~\ref{fig:nbobservables-baseline} compares full-$M$ and v3 no-$M$+$M_S$
at both fillings. Figures~\ref{fig:nbobservables13} and
\ref{fig:nbobservables23} display all three externally projected sectors. At $\nu=1/3$, the first row is the
training-energy minimum $P_0$ and is the selected out-$C_3$ result; $P_1$ and
$P_2$ are retained only for a controlled sector comparison. At $\nu=2/3$,
$P_2$ is the corresponding minimum. The panels are regenerated from the
saved final training ensembles, with no fresh validation chains. Each contains 4,128 configurations;
energy selection and observable estimation use the same production-chain
convention as the optimization protocol of Ref.~\cite{luo}.

\subsection{Full-ratio 1-RDM and fixed bare-band projector}

For $a=(n,\bm k)$ and $b=(m,\bm q)$, the estimator is
\begin{equation}
\begin{split}
 \rho_{ab}&=\langle c_b^\dagger c_a\rangle\\
 &=2N_e\,\mathbb E_{X\sim|\Psi|^2,i,x'}\!\left[
 (\phi_a^0(x_i))^*\phi_b^0(x')
 \left(\frac{\Psi(X_{i\to x'})}{\Psi(X)}\right)^*\right].
\end{split}
\label{eq:rdmratio}
\end{equation}
The factor two accounts for uniform auxiliary sampling of the two layers.  In
every ratio, replacing one particle triggers a complete new ansatz call: the
code rebuilds the graph and recomputes all particles' $\delta\bm r_j$,
$J_{\bm k j}$, and either the $M Q$ determinant or all $M_S$-weighted sector
minors. The
projector is instead the fixed bare function
\begin{equation}
 \phi^0_{n\bm k}(\bm r,l)=\sum_{\bm G}u^l_{n\bm k\bm G}
 e^{i(\bm k+\bm G)\cdot\bm r};
\end{equation}
no backflow phase is multiplied into it.  Thus neural dressing appears once,
only through the complete wavefunction ratio.

All 4,128 final training configurations and 16 auxiliary points per configuration
are used to estimate the $45\times45$ five-band-projected matrix.  The band
populations are
\begin{equation}
 W_n=\frac1{N_e}\sum_{\bm k}\rho_{n\bm k,n\bm k}.
\end{equation}
\begin{table}[H]
\centering
\caption{Five-band-projected 1-RDM results from final training samples.
Band weights are percentages; parentheses are eight-block estimates.  ED
weights use the exact five-band one-body density matrix.}
\label{tab:adaptive-rdm}
\resizebox{\columnwidth}{!}{%
\begin{tabular}{@{}llrrrrrrr@{}}
\toprule
$\nu$ & state & trace & $W_1$ & $W_2$ & $W_3$ & $W_4$ & $W_5$ & min eig.\\
\midrule
$1/3$ & five-band ED & $3$ & 96.47 & 2.95 & 0.52 & 0.05 & 0.01 & ---\\
$1/3$ & no $M$+$M_S$ & $3.167(78)$ & $95.41(1.21)$ & $5.94(64)$ & $2.67(1.25)$ & $1.01(1.57)$ & $0.54(48)$ & $-0.185$\\
$1/3$ & outer $P_0$+$M_S$ & $3.013(70)$ & $96.47(93)$ & $0.47(1.16)$ & $1.42(1.11)$ & $0.83(75)$ & $1.23(1.02)$ & $-0.111$\\
$1/3$ & outer $P_1$+$M_S$ & $3.030(90)$ & $81.40(1.05)$ & $17.73(1.38)$ & $2.15(93)$ & $1.04(1.40)$ & $-1.34(79)$ & $-0.103$\\
$1/3$ & outer $P_2$+$M_S$ & $2.961(72)$ & $94.65(1.17)$ & $2.35(74)$ & $0.53(94)$ & $0.77(59)$ & $0.39(36)$ & $-0.121$\\
\midrule
$2/3$ & five-band ED & $6$ & 88.95 & 7.59 & 2.53 & 0.76 & 0.16 & ---\\
$2/3$ & no $M$+$M_S$ & $6.087(158)$ & $91.29(1.25)$ & $7.30(93)$ & $2.47(73)$ & $0.27(63)$ & $0.13(77)$ & $-0.151$\\
$2/3$ & outer $P_0$+$M_S$ & $5.744(148)$ & $87.27(73)$ & $6.31(1.24)$ & $1.35(68)$ & $1.43(1.09)$ & $-0.62(92)$ & $-0.288$\\
$2/3$ & outer $P_1$+$M_S$ & $5.959(174)$ & $87.10(69)$ & $8.57(1.16)$ & $3.03(1.13)$ & $1.58(71)$ & $-0.96(1.35)$ & $-0.240$\\
$2/3$ & outer $P_2$+$M_S$ & $6.038(170)$ & $87.87(1.86)$ & $6.48(96)$ & $3.02(59)$ & $2.49(1.01)$ & $0.78(85)$ & $-0.205$\\
\bottomrule
\end{tabular}%
}
\end{table}
Most projected traces are compatible with $N_e$ at the stated precision, but
the matrices are only Hermitized, not projected onto the positive-semidefinite
cone. Their negative minimum eigenvalues therefore forbid a
stable natural-orbital eigenspace interpretation.  The diagonal band
populations are more robust and reveal a real sector dependence.  At
$\nu=1/3$, $P_0$ and $P_2$ retain 96.47(93)\% and 94.65(1.17)\% band-1
weight, whereas $P_1$ has $81.40(1.05)\%$ and $17.73(1.38)\%$ in band 2; five-band ED has
$W_1=96.47\%$ and $W_2=2.95\%$.  At $\nu=2/3$, all three sectors are closer
to the ED distribution $(88.95,7.59,2.53,0.76,0.16)\%$; the lowest-training
$P_2$ gives $(87.87,6.48,3.02,2.49,0.78)\%$. Small negative entries remain
consistent with auxiliary-coordinate noise.

The estimator now caches $\Psi_m(X)$ once per configuration block and jointly
batches all 16 replacement coordinates.  Each replacement still rebuilds the
complete graph, all $\delta\bm r_j$, $J$, the $M_S$ determinant mixture,
the projected three-branch sum and the wavefunction ratio.  On the A30 GPUs,
the 4,128-by-16 outer-sector measurement processes take 4.1--4.2 seconds at
$\nu=1/3$ and 5.9--6.2 seconds at $\nu=2/3$.
Checkpoint conversion plus all three measurements takes 21.7 and 25.6 seconds,
respectively,
on the scheduled dual-GPU critical path. Training and postprocessing are timed
separately in Table~\ref{tab:ed-qns-efficiency}.
The focused bare-Slater regression sets $\delta\bm r=0$, $J=1$, fixes a
complex $2\times9$ matrix $M$ with orthonormal rows, and samples the resulting
single determinant.  In a complete nine-orbital test basis the exact matrix is
$\rho=M^{\mathsf T}M^*$, with trace two.  With 4,096 walkers, 100 burn sweeps
and 32 auxiliary draws, the measured trace is $2.01047\pm0.01459$, the
Frobenius error is 0.05272, the Hermiticity residual is zero, and the minimum
sampled eigenvalue is $-0.02072$.  This unit test verifies occupation,
Hermiticity, PSD within sampling noise and full-basis trace normalization.

\subsection{Full structure factor, connected part, and density order}

The structure factor is evaluated directly as a two-body observable, not
reconstructed from the 1-RDM:
\begin{align}
 S_{\rm full}(\bm q)&=\frac{\langle|\rho_{\bm q}|^2\rangle}{N_e},\\
 S_c(\bm q)&=\frac{\langle|\rho_{\bm q}|^2\rangle-
 |\langle\rho_{\bm q}\rangle|^2}{N_e},\\
 m_{\bm q}&=\frac{|\langle\rho_{\bm q}\rangle|}{N_e}.
\end{align}
The ED comparison uses $S_{\rm full}$ with the same normalization.  At
$\nu=1/3$, outer $P_0,P_1,P_2$ give maximum nonzero-momentum
$S_{\rm full}=1.5058(143),1.1281(145),1.2294(102)$, respectively, compared
with the five-band ED value 1.4381.  At $\nu=2/3$ the corresponding values are
$0.7367(78),0.6799(101),0.7123(99)$, compared with ED 0.6197.  The sector
dependence is therefore visible in both the 1-RDM and the genuine two-body
observable; it is not a plotting artifact.

For $m=0,1,2$, the largest nonzero-$q$ density amplitudes are
$0.00728(293),0.00702(229),0.00856(352)$ at $\nu=1/3$ and
$0.00365(259),0.00261(166),0.00620(266)$ at $\nu=2/3$. These small values
show negligible static translation-breaking expectation at the sampled
reciprocal vectors; the large $S_{\rm full}$ peaks are density fluctuations.
The real-space panels in Figs.~\ref{fig:nbobservables13} and
\ref{fig:nbobservables23} include the intra-moir\'e-cell Bloch density and
should not be identified with a nonzero CDW order parameter by visual
contrast alone.

\begin{figure*}[t]
\centering
\includegraphics[width=.98\textwidth]{fig7_neural_bloch_observables.pdf}
\caption{\textbf{Full-$M$ and v3 no-$M$+$M_S$ observables.} The first two
rows are $\nu=1/3$ and the last two are $\nu=2/3$; at each filling full $M$
is followed by no $M$+$M_S$. Columns show bare-band-1 momentum occupation,
the full static structure factor, and folded charge density. The v3 no-$M$
panels use the newly evaluated 4,128-configuration final ensembles.}
\label{fig:nbobservables-baseline}
\end{figure*}

\begin{figure*}[t]
\centering
\includegraphics[width=.98\textwidth]{fig7_neural_bloch_observables_nu1of3.pdf}
\caption{\textbf{$\nu=1/3$ outer-sector observables from final
training samples.} Rows are complete-wavefunction projectors $P_0,P_1,P_2$;
the first row, $P_0$, is explicitly marked as the selected outer-$C_3$ state because it has the
lowest final-ten training energy.
Columns show bare-band-1 momentum occupation from the complete-ratio 1-RDM,
the full static structure factor with ED normalization on the complete
hexagonal BZ, and folded charge density over a two-by-two moir\'e-cell
parallelogram. Shared color limits expose the strong $P_1$ band-mixing and
structure-factor difference without rescaling individual sectors.}
\label{fig:nbobservables13}
\end{figure*}

\begin{figure*}[t]
\centering
\includegraphics[width=.98\textwidth]{fig7_neural_bloch_observables_nu2of3.pdf}
\caption{\textbf{$\nu=2/3$ outer-sector observables from final
training samples.} Rows are $P_0,P_1,P_2$ and columns use the same definitions
and global color limits as Fig.~\ref{fig:nbobservables13}.  All panels use
4,128 configurations and 16 jointly batched auxiliary replacements per
configuration.  The three sectors are shown separately because both their
band populations and two-body structure factors differ beyond a simple
identical-map interpretation. The $P_2$ row is marked as the lowest
final-ten-energy sector.}
\label{fig:nbobservables23}
\end{figure*}

\section{Conclusion}

The implemented baseline follows the printed Luo--Fu Appendix-B structure:
fixed mean-field band-1 plane-wave coefficients, periodic node and edge
features, trainable random auxiliary features, two unshared width-32 message
iterations with the literal self-inclusive $\sum_jm_{ij}$, coordinate
backflow, direct complex $J_{\bm k i}$ applied exactly once, and one final
generalized Slater construction. Full $M$ mixes momenta at the single-particle
level and therefore generally mixes total-momentum sectors. Version 3 no-$M$
instead uses $M_S\in\mathbb C^{1\times12}$ to combine all and only the
selected-momentum determinants in the exact total-$\Gamma$ sector.

The outer construction leaves that raw network unrestricted and projects the
complete $M_S$ determinant mixture with $P_m$. Character and
projector-algebra residuals are below $6.4\times10^{-14}$ and
$7.6\times10^{-15}$, all three projected norms are nonzero, and local energy,
Metropolis ratios and SR derivatives all use the projected state.  It adds no
trainable parameter to the 78,124-parameter v3 no-$M$ model.

On the final-ten training tail, v3 no-$M$ reaches $-37.3037$ meV per particle
at $\nu=1/3$, while outer $P_0$ is the lowest of the three projected sectors
at $-37.3436$ meV under the same v3.2 shared-pool protocol.
Its $P_1$ and $P_2$ tails still descend. At $\nu=2/3$, v3 no-$M$ reaches $-52.7666$ meV and the
v3.2 same-seed, shared-pool run gives $P_2$ as the lowest projected sector
at $-52.8838$ meV. P0 and P1 reach $-52.5922$ and $-51.5740$ meV per
particle. P1 still descends near the end; this single-initialization experiment
does not establish convergence or resolve the seed-sensitivity question.

The all-sector observables show that $m=0,1,2$ cannot be represented by one
duplicated map: at $\nu=1/3$, $P_1$ has markedly enhanced band-2 weight and a
smaller structure-factor peak, while at $\nu=2/3$ the differences are weaker
but remain resolved.  Thus the complete-state projector changes both energy
and correlations, rather than merely relabeling the same generalized
determinant.
The complete-ratio 1-RDM and direct two-body structure factors remain useful
diagnostics, but their non-PSD noise and the single-seed finite optimization
preclude a natural-orbital or topological claim. The next controlled tests are
longer multi-seed optimization and sector-selection replication.
